% ==============================================================================
% تقرير المشروع النهائي لمقرر تصميم وبناء المترجمات (Compilers Design)
% مشروع مترجم ومحرر لغة بيان العربية (Bayan Compiler & Arabic VS Code IDE)
% الرمز المعتمد للتسليم: RPT-G01
% تم الإعداد للتجميع عبر XeLaTeX أو Overleaf
% ==============================================================================

\documentclass[12pt,a4paper]{report}

% حزم دعم اللغة العربية والخطوط (XeLaTeX / LuaLaTeX)
\usepackage{fontspec}
\usepackage{polyglossia}
\setmainlanguage[numerals=maghrib]{arabic}
\setotherlanguage{english}

% تعيين خطوط احترافية تدعم العربية والإنجليزية
\defaultfontfeatures{Ligatures=TeX}
% في حال التجميع على Overleaf أو بيئة TeX Live القياسية يتوفر خط Amiri
\newfontfamily\arabicfont[Script=Arabic]{Amiri}
\newfontfamily\arabicfontsf[Script=Arabic]{Amiri}
\newfontfamily\englishfont{DejaVu Sans}

% الحزم الرياضية والرسومية وتنسيق الصفحات
\usepackage{amsmath,amssymb}
\usepackage{graphicx}
\usepackage{geometry}
\geometry{top=2.5cm, bottom=2.5cm, left=2.5cm, right=2.5cm}
\usepackage{fancyhdr}
\usepackage{titlesec}
\usepackage{xcolor}
\usepackage{tcolorbox}
\usepackage{listings}
\usepackage{tabularx}
\usepackage{booktabs}
\usepackage{hyperref}
\usepackage{caption}
\usepackage{float}

% إعدادات الروابط التشعبية
\hypersetup{
    colorlinks=true,
    linkcolor=blue!70!black,
    citecolor=green!60!black,
    urlcolor=blue!80!black
}

% تعريف لوحة ألوان Visual Studio Code
\definecolor{vscodeBlue}{RGB}{0,122,204}
\definecolor{vscodeGreen}{RGB}{78,201,176}
\definecolor{vscodeDark}{RGB}{30,30,30}
\definecolor{vscodeGray}{RGB}{240,242,245}
\definecolor{codeKw}{RGB}{86,156,214}
\definecolor{codeStr}{RGB}{206,145,120}
\definecolor{codeComment}{RGB}{106,153,85}

% إعدادات قوائم الأكواد المصدرية (Listings)
\lstset{
    basicstyle=\small\ttfamily\englishfont,
    numbers=left,
    numberstyle=\tiny\color{gray},
    stepnumber=1,
    numbersep=8pt,
    backgroundcolor=\color{vscodeGray},
    showspaces=false,
    showstringspaces=false,
    showtabs=false,
    frame=single,
    rulecolor=\color{gray!30},
    tabsize=4,
    captionpos=b,
    breaklines=true,
    breakatwhitespace=false,
    keywordstyle=\color{codeKw}\bfseries,
    commentstyle=\color{codeComment}\itshape,
    stringstyle=\color{codeStr}
}

% تنسيق الترويسة والتذييل
\pagestyle{fancy}
\fancyhf{}
\fancyhead[R]{\small\color{gray}مشروع مترجم ومحرر لغة بيان العربية}
\fancyhead[L]{\small\color{gray}مقرر تصميم وبناء المترجمات 2025/2026}
\fancyfoot[C]{\thepage}
\renewcommand{\headrulewidth}{0.4pt}
\renewcommand{\footrulewidth}{0.4pt}

\begin{document}

% ==============================================================================
% الصفحة الأولى: صفحة الغلاف الرسمية (Cover Page)
% ==============================================================================
\begin{titlepage}
    \centering
    \vspace*{-0.5cm}
    
    \begin{minipage}[c]{0.35\textwidth}
        \flushright
        {\bfseries الجمهورية اليمنية}\\[0.15cm]
        {\small وزارة التعليم العالي والبحث العلمي}\\[0.15cm]
        {\bfseries جامعة إب}\\[0.15cm]
        {\small كلية الحاسوب وتكنولوجيا المعلومات}\\[0.15cm]
        {\small قسم علوم الحاسوب}
    \end{minipage}
    \hfill
    \begin{minipage}[c]{0.26\textwidth}
        \centering
        \includegraphics[width=2.8cm]{images/ibb_university_logo.png}
    \end{minipage}
    \hfill
    \begin{minipage}[c]{0.35\textwidth}
        \flushleft
        \begin{english}
        {\bfseries Republic of Yemen}\\[0.15cm]
        {\scriptsize Ministry of Higher Education}\\[0.15cm]
        {\bfseries Ibb University}\\[0.15cm]
        {\scriptsize Faculty of Computer \& IT}\\[0.15cm]
        {\scriptsize Computer Science Dept.}
        \end{english}
    \end{minipage}
    
    \vspace{0.4cm}
    \hrule height 1.5pt
    \vspace{0.08cm}
    \hrule height 0.5pt
    \vspace{0.8cm}
    
    {\large\color{vscodeBlue}\bfseries تقرير المشروع النهائي لنيل درجة متطلب مقرر:}\\[0.3cm]
    {\LARGE\bfseries تصميم وبناء المترجمات}\\[0.2cm]
    {\large\textenglish{(Compilers Design)}}\\[0.8cm]
    
    \begin{tcolorbox}[colback=vscodeBlue!5!white,colframe=vscodeBlue!80!black,arc=2mm,boxrule=1.2pt]
        \centering
        \vspace{0.2cm}
        {\huge\bfseries مشروع مترجم ومحرر لغة بيان}\\[0.3cm]
        {\Large\bfseries\textenglish{Bayan Arabic Compiler \& Arabic IDE Studio}}\\[0.25cm]
        {\normalsize\color{vscodeBlue!90!black}\bfseries مشروع تطبيقي متكامل لتصميم المترجمات وتوليد البرامج التنفيذية المستقلة (Windows PE EXE)}\\[0.2cm]
    \end{tcolorbox}
    
    \vspace{0.8cm}
    
    \begin{minipage}[t]{0.45\textwidth}
        \begin{tcolorbox}[colback=white,colframe=gray!60,arc=2mm,title=\textbf{إشراف الدكتور الفاضل:}]
            \centering
            \vspace{0.2cm}
            {\Large\bfseries د. خالد الكحسة}\\[0.2cm]
            {\small أستاذ مقرر تصميم وبناء المترجمات}\\[0.2cm]
        \end{tcolorbox}
    \end{minipage}
    \hfill
    \begin{minipage}[t]{0.50\textwidth}
        \begin{tcolorbox}[colback=white,colframe=gray!60,arc=2mm,title=\textbf{إعداد طلاب المجموعة (رقم: 01):}]
            \small
            \begin{tabular}{r l}
                \textbf{1. محمد رشاد الإدريسي} & \textenglish{(M. Rashad Al-Edrisi)} \\
                \textbf{2. أسامة شميس} & \textenglish{(Osama Shomis)} \\
                \textbf{3. أسامة العقاب} & \textenglish{(Osama Al-Oqab)} \\
                \textbf{4. محمد فيصل الدعيس} & \textenglish{(M. Faisal Al-Doais)} \\
                \textbf{5. شعيب جازم} & \textenglish{(Shoaib Jazem)} \\
            \end{tabular}
        \end{tcolorbox}
    \end{minipage}
    
    \vfill
    
    \hrule height 0.5pt
    \vspace{0.3cm}
    {\normalsize\bfseries الفصل الدراسي الثاني — العام الجامعي: 2025 / 2026 م — 1447 هـ}\\[0.2cm]
    
\end{titlepage}

% ==============================================================================
% فهرس المحتويات
% ==============================================================================
\tableofcontents
\newpage

% ==============================================================================
% الصفحة الثانية: وصف المشروع والبيئة التقنية والمكتبات
% ==============================================================================
\chapter{وصف المشروع والبيئة التقنية والمعمارية}

\section{وصف المشروع (Project Overview)}
مشروع **لغة بيان (Bayan)** هو بيئة برمجية متكاملة تهدف إلى توفير لغة برمجة عربية عالية المستوى، مدعومة بمترجم مستقل مبني من الصفر (From Scratch)، ومحرر برمجي تفاعلي حديث يحاكي واجهة وتجربة المستخدم لبرنامج **Visual Studio Code**، مع دعم أصيل وكامل لاتجاه الكتابة العربي من اليمين إلى اليسار (RTL).

ينقسم المشروع وفق وثيقة المتطلبات إلى نظامين مستقلين تماماً:
\begin{enumerate}
    \item \textbf{محرك المترجم المستقل (BayanCompiler.exe):} تطبيق سطر أوامر (CLI) ينفذ دورة الترجمة الكاملة المكونة من 9 مراحل متتالية؛ تبدأ من قراءة الكود المصدري العربي، وتمر بالتحليل المعجمي، والتحليل النحوي التنبؤي الصارم $LL(1)$، وبناء شجرة الإعراب المجردة (AST)، والتحليل الدلالي مع إدارة مجالات الرؤية في جدول الرموز، ثم توليد الشفرة الوسيطة ثلاثية العناوين (TAC)، وتوليد شفرة التجميع لمعالجات $x86$، وتوليد لغة وسيطة $CIL/MSIL$ واستدعاء مجمّع مايكروسوفت $ilasm.exe$ لإنتاج ملف تنفيذي أصلي $EXE$ يعمل مباشرة في بيئة ويندوز.
    \item \textbf{محرر اللغة وبيئة التطوير (Bayan.Editor.WinForms.exe):} بيئة تطوير مرئية متطورة توفر تجربة مبرمج احترافية؛ تتيح إدارة الملفات المتعددة، تمييز الكلمات المفتاحية باللون، استكشاف أمثلة اللغة، وتفحص مخرجات مراحل الترجمة التسع في تبويبات مخصصة، مع طرفية تفاعلية لتشغيل البرامج وقائمة أخطاء تقفز مباشرة إلى السطر والعمود.
\end{enumerate}

\section{العلاقة المعمارية المنفصلة (Decoupled Architecture)}
لا يعتمد المحرر على استدعاء دوال المترجم داخلياً، بل يحقق الفصل التام (Decoupling) عبر تشغيل المترجم كعملية خارجية منفصلة (Inter-Process Communication):
\begin{center}
    \textbf{المستخدم $\longleftrightarrow$ محرر بيان (IDE) $\longleftrightarrow$ استدعاء سطر الأوامر (CLI) $\longleftrightarrow$ محرك المترجم (BayanCompiler) $\longleftrightarrow$ الملف التنفيذي (EXE)}
\end{center}

\section{لغات البرمجة المستخدمة والإصدارات (Programming Languages)}
\begin{itemize}
    \item \textbf{C\# (الإصدار 10.0):} تم بناء كامل المحرك والمحرر بلغة C\# الحديثة مع الاستفادة من ميزات التنميط القوي، ومطابقة الأنماط (Pattern Matching)، وميزات إدارة الذاكرة الآمنة.
    \item \textbf{منصة التشغيل:} \textenglish{.NET 6.0 LTS (Long Term Support)} إصدار \textenglish{6.0.100}.
\end{itemize}

\section{الأدوات والمكتبات الجاهزة (Tools \& Libraries)}
وفقاً للمعيار الخامس في وثيقة متطلبات المشروع، يوضح الجدول التالي كافة الأدوات والمكتبات المستخدمة ودور كل منها:

\begin{table}[H]
    \centering
    \caption{الأدوات والمكتبات الجاهزة المستخدمة في المشروع}
    \begin{tabularx}{\textwidth}{l p{3.2cm} l X}
        \toprule
        \textbf{اسم الأداة / المكتبة} & \textbf{الإصدار} & \textbf{المرحلة} & \textbf{الغرض وتفاصيل الاستخدام} \\
        \midrule
        \textenglish{Microsoft IL Assembler (ilasm.exe)} & \textenglish{4.8.9037.0} & توليد الملف التنفيذي & تجميع كود CIL/MSIL وتحويله إلى ملف تنفيذي أصلي بنظام Windows PE (32/64-bit Console). \\
        \addlinespace
        \textenglish{System.Text.Json} & \textenglish{6.0.0} & النقل والتصدير & تسلسل شجرة الإعراب (AST) وجدول الرموز (Symbol Table) إلى نسق JSON منظم. \\
        \addlinespace
        \textenglish{Windows Forms} & \textenglish{6.0.0 (.NET 6)} & واجهة المحرر & بناء واجهة VS Code العربية الداعمة لـ RTL وعناصر التحكم الرسومية والطرفية. \\
        \addlinespace
        \textenglish{System.Drawing.Common} & \textenglish{6.0.0} & الرسوم والأيقونات & رسم شريط الأنشطة وحساب ألوان السمة الداكنة الداعمة لمعايير VS Code. \\
        \bottomrule
    \end{tabularx}
\end{table}

\newpage

% ==============================================================================
% من الصفحة الثالثة: شرح وتوثيق مراحل المترجم وأكواد المشروع
% ==============================================================================
\chapter{المراحل التفصيلية للمترجم وشرح الأكواد المصدرية (وفق البنود الـ 9)}

امتثالاً للبند الثامن من وثيقة تعليمات المشروع، يقدم هذا الفصل شرحاً تحليلياً شاملاً لكافة أجزاء ومراحل المترجم ومحرر اللغة، يغطي لكل جزء النقاط التسع الإلزامية المحددة في معايير التقييم مع مقتطفات برمجية حقيقية من صلب المشروع:

% ------------------------------------------------------------------------------
\section{المحلل المعجمي واستخراج الرموز (Lexer.cs)}

\begin{tcolorbox}[colback=white,colframe=vscodeBlue!80!black,arc=2mm,title=\textbf{بطاقة المواصفات البرمجية للمحلل المعجمي (Lexical Analyzer Specification)}]
\begin{tabularx}{\linewidth}{p{4.2cm} X}
    \textbf{1. وظيفة الجزء من الكود} & قراءة الكود المصدري العربي حرفاً بحرف وتحليله معجمياً، وتفكيكه إلى تدفق من الرموز المعجمية (Tokens)، مع عزل التعليقات والفراغات، واحتساب إحداثيات السطر والعمود بدقة، وتشخيص الرموز غير المعرفة. \\
    \midrule
    \textbf{2. المدخلات (Inputs)} & النص البرمجي المصدري المكتوب بلغة بيان بترميز \textenglish{UTF-8} الممرر من المحرر أو سطر الأوامر. \\
    \midrule
    \textbf{3. المخرجات (Outputs)} & كائن \textenglish{LexResult} يحتوي على قائمة الرموز المعجمية المتسلسلة \textenglish{List\textless Token\textgreater}، بالإضافة إلى سجل الأخطاء المعجمية \textenglish{DiagnosticBag} إن وُجدت. \\
    \midrule
    \textbf{4. المتغيرات المستخدمة} & \textenglish{\_position} (مؤشر موقع القراءة الحالي في النص)، \textenglish{\_line} (رقم السطر النشط يبدأ من 1)، \textenglish{\_column} (رقم العمود النشط)، \textenglish{\_startPosition / \_startLine / \_startColumn} (إحداثيات بداية الرمز)، \textenglish{c} (المحرف الحالي المفحوص). \\
    \midrule
    \textbf{5. هياكل البيانات المستخدمة} & \textenglish{enum TokenType} (يحوي أكثر من 40 نوعاً للرموز)، \textenglish{class Token} (سجل يضم النوع، القيمة النصية Lexeme، السطر، والعمود)، قاموس الكلمات المحجوزة \textenglish{Dictionary\textless string, TokenType\textgreater} لمطابقة الكلمات المفتاحية في زمن ثابت $O(1)$. \\
    \midrule
    \textbf{6. الدوال أو الإجراءات المستخدمة} & \textenglish{Lex()} (الدالة الرئيسية لبدء المسح)، \textenglish{NextToken()} (استخراج الرمز التالي)، \textenglish{ReadIdentifierOrKeyword()} (قراءة المعرفات والكلمات العربية)، \textenglish{ReadNumber()} (معالجة الأعداد الصحيحة والحقيقية)، \textenglish{ReadString()} (قراءة السلاسل النصية)، \textenglish{Peek()} (استشراف المحرف التالي). \\
    \midrule
    \textbf{7. المكتبات المستخدمة} & \textenglish{System, System.Collections.Generic, System.Text} ومكتبة الواجهة الأمامية للمشروع \textenglish{BayanCompiler.Frontend}. \\
    \midrule
    \textbf{8. الملفات المرتبطة بالجزء} & \textenglish{Token.cs, TokenType.cs, SourceText.cs, Diagnostic.cs, DiagnosticBag.cs}. \\
    \midrule
    \textbf{9. كيفية ارتباطه بباقي المشروع} & يمثل المرحلة الأولى (بوابة خط الأنابيب)؛ يستقبل النص المباشر من واجهة المحرر، ويمرر مصفوفة الرموز (Tokens) مباشرة كمدخل أساسي للمحلل النحوي التنبؤي \textenglish{DoctorLl1Parser.cs}. \\
\end{tabularx}
\end{tcolorbox}

\begin{lstlisting}[language={[Sharp]C}, caption={مقتطف حقيقي من المحلل المعجمي Lexer.cs}]
public LexResult Lex()
{
    var tokens = new List<Token>();
    while (_position < _text.Length)
    {
        char current = Current;
        if (char.IsWhiteSpace(current)) { Advance(); continue; }
        if (current == '/' && Peek(1) == '/') { SkipLineComment(); continue; }

        if (char.IsLetter(current) || current == '_')
        {
            tokens.Add(ReadIdentifierOrKeyword());
            continue;
        }
        if (char.IsDigit(current))
        {
            tokens.Add(ReadNumber());
            continue;
        }
        if (current == '"')
        {
            tokens.Add(ReadString());
            continue;
        }
        tokens.Add(ReadOperatorOrDelimiter());
    }
    tokens.Add(new Token(TokenType.EndOfFile, "\0", new SourceSpan(_position, 0, _line, _column)));
    return new LexResult(tokens, _diagnostics.ToReadOnlyList());
}
\end{lstlisting}

\newpage

% ------------------------------------------------------------------------------
\section{المحلل النحوي التنبؤي LL(1) وبناء الشجرة المجردة (DoctorLl1Parser.cs \& DoctorAstBuilder.cs)}

\begin{tcolorbox}[colback=white,colframe=vscodeBlue!80!black,arc=2mm,title=\textbf{بطاقة المواصفات البرمجية للمحلل النحوي (Predictive LL(1) Parser Specification)}]
\begin{tabularx}{\linewidth}{p{4.2cm} X}
    \textbf{1. وظيفة الجزء من الكود} & التحقق من صحة القواعد النحوية للكود المصدري وفق قواعد الطبيب عبر خوارزمية التحليل التنبؤي الصارم $LL(1)$ غير العودي والخالي من التراجع، مع بناء شجرة الاشتقاق (Parse Tree) وشجرة الإعراب المجردة (AST). \\
    \midrule
    \textbf{2. المدخلات (Inputs)} & مصفوفة الرموز المعجمية \textenglish{IReadOnlyList\textless Token\textgreater} الناتجة من المحلل المعجمي. \\
    \midrule
    \textbf{3. المخرجات (Outputs)} & كائن \textenglish{DoctorParseResult} يتضمن جذر الشجرة المجردة \textenglish{ProgramNode Program}، شجرة الاشتقاق \textenglish{ParseTreeNode ParseTree}، سجل تتبع المكدس \textenglish{Ll1Trace}، وجدول التحليل \textenglish{Ll1Table}. \\
    \midrule
    \textbf{4. المتغيرات المستخدمة} & \textenglish{\_tokens} (قائمة الرموز الممررة)، \textenglish{\_position} (مؤشر الرمز المفحوص)، \textenglish{\_parseStack} (مكدس التحليل التنبؤي للرموز غير النهائية والنهائية)، \textenglish{\_traceBuilder} (سجل توثيق خطوات التحليل). \\
    \midrule
    \textbf{5. هياكل البيانات المستخدمة} & \textenglish{Stack\textless string\textgreater} (مكدس رموز القواعد)، جدول توجيه التحليل $M[A, a]$ ممثلاً بـ \textenglish{Dictionary\textless(string, TokenType), List\textless string\textgreater\textgreater}، وشجرة كائنية هرمية للعقد: \textenglish{ProgramNode, VariableDeclarationNode, AssignmentNode, IfNode, LoopNode, ProcedureNode}. \\
    \midrule
    \textbf{6. الدوال أو الإجراءات المستخدمة} & \textenglish{Parse()} (بدء التحليل النحوي)، \textenglish{Predict()} (استخراج قاعدة الإنتاج المطابقة من جدول LL1)، \textenglish{Match()} (مطابقة الرمز النهائي واستهلاكه)، \textenglish{BuildAst()} (تحويل شجرة الاشتقاق لعقد AST)، \textenglish{FormatTrace()} (تنسيق خطوات التحليل). \\
    \midrule
    \textbf{7. المكتبات المستخدمة} & \textenglish{System, System.Collections.Generic, BayanCompiler.Syntax, BayanCompiler.Lexing}. \\
    \midrule
    \textbf{8. الملفات المرتبطة بالجزء} & \textenglish{DoctorAstBuilder.cs, DoctorSpecificationNodes.cs, ParseResult.cs, AstNode.cs, Statements.cs, Expressions.cs}. \\
    \midrule
    \textbf{9. كيفية ارتباطه بباقي المشروع} & يستهلك رموز المحلل المعجمي، ويبني شجرة الإعراب المجردة (AST) التي تمثل الهيكل المحوري الممرر للمحلل الدلالي \textenglish{DoctorSemanticAnalyzer.cs}. \\
\end{tabularx}
\end{tcolorbox}

\begin{lstlisting}[language={[Sharp]C}, caption={مقتطف حقيقي من التحليل النحوي التنبؤي DoctorLl1Parser.cs}]
public DoctorParseResult Parse()
{
    var parseStack = new Stack<string>();
    parseStack.Push(DoctorGrammar.RuleEndMarker);
    parseStack.Push(DoctorGrammar.StartSymbol); // "برنامج"

    while (parseStack.Count > 0)
    {
        string top = parseStack.Pop();
        Token currentToken = Current;

        if (IsTerminal(top))
        {
            if (MatchTerminal(top, currentToken))
            {
                Advance(); // استهلاك الرمز بعد التاكد من صحته
            }
            else
            {
                ReportSyntaxError(top, currentToken);
                return DoctorParseResult.Failure(_diagnostics);
            }
        }
        else // رمز غير نهائي: تطبيق قاعدة الانتاج من جدول LL(1)
        {
            var production = LookupProduction(top, currentToken.Type);
            for (int i = production.Count - 1; i >= 0; i--)
            {
                if (production[i] != DoctorGrammar.Epsilon)
                    parseStack.Push(production[i]);
            }
        }
    }
    return DoctorParseResult.Success(_astBuilder.Build());
}
\end{lstlisting}

\newpage

% ------------------------------------------------------------------------------
\section{التحليل الدلالي وإدارة مجالات الرؤية وجدول الرموز (DoctorSemanticAnalyzer.cs \& SymbolTable.cs)}

\begin{tcolorbox}[colback=white,colframe=vscodeBlue!80!black,arc=2mm,title=\textbf{بطاقة المواصفات البرمجية للتحليل الدلالي (Semantic Analyzer Specification)}]
\begin{tabularx}{\linewidth}{p{4.2cm} X}
    \textbf{1. وظيفة الجزء من الكود} & التحقق من صحة المعنى والمنطق البرمجي (Semantic Rules)، فرض التصريح المسبق عن المتغيرات، حماية الثوابت من التعديل، التحقق من تطابق وتوافق الأنواع في العمليات الرياضية والمقارنات، وإدارة مجالات الرؤية المتداخلة (Scopes). \\
    \midrule
    \textbf{2. المدخلات (Inputs)} & شجرة الإعراب المجردة \textenglish{ProgramNode} الناتجة من مرحلة التحليل النحوي. \\
    \midrule
    \textbf{3. المخرجات (Outputs)} & كائن \textenglish{DoctorSemanticResult} يحتوي على حالة النجاح \textenglish{bool Succeeded}، جدول الرموز الهرمي \textenglish{SymbolTable} شاملاً كافة المعرفات وأنواعها ومجالاتها، ومصفوفة التشخيصات الدلالية \textenglish{Diagnostics}. \\
    \midrule
    \textbf{4. المتغيرات المستخدمة} & \textenglish{\_currentScope} (النطاق النشط حالياً أثناء التجوال)، \textenglish{\_globalScope} (نطاق البرنامج العام)، \textenglish{\_diagnostics} (سجل تجميع الأخطاء الدلالية)، \textenglish{expectedType} (النوع المتوقع للطرفين)، \textenglish{actualType} (النوع المشتق للتعبير). \\
    \midrule
    \textbf{5. هياكل البيانات المستخدمة} & كلاس \textenglish{SymbolTable} بآلية الإشارة للنطاق الأب \textenglish{ParentScope}، قاموس الرموز \textenglish{Dictionary\textless string, Symbol\textgreater}، كلاس \textenglish{Symbol} (يحمل الاسم، الفئة، نوع البيانات، هل هو للقراءة فقط ReadOnly، ورقم السطر)، وتعداد أنواع اللغة \textenglish{LanguageType} (صحيح، حقيقي، منطقي، حرفي، خيط رمزي، قائمة، سجل، إجراء). \\
    \midrule
    \textbf{6. الدوال أو الإجراءات المستخدمة} & \textenglish{Analyze()} (نقطة الدخول للتحليل)، \textenglish{VisitStatement()} (فحص الجمل والتصريحات)، \textenglish{VisitExpression()} (اشتقاق نوع التعبير)، \textenglish{Declare()} (تسجيل رمز جديد بالنطاق)، \textenglish{Lookup()} (البحث عن رمز صعوداً عبر النطاقات)، \textenglish{EnterScope() / ExitScope()} (إدارة بيئات الرؤية). \\
    \midrule
    \textbf{7. المكتبات المستخدمة} & \textenglish{System, System.Collections.Generic, BayanCompiler.Syntax, BayanCompiler.Frontend}. \\
    \midrule
    \textbf{8. الملفات المرتبطة بالجزء} & \textenglish{Symbol.cs, SymbolTable.cs, LanguageType.cs, TypeDescriptor.cs, DoctorSemanticAnalyzer.cs}. \\
    \midrule
    \textbf{9. كيفية ارتباطه بباقي المشروع} & يستقبل شجرة AST من المحلل النحوي، ويؤكد سلامتها من أخطاء الأنواع والمجالات، ويزود مولد الشفرة الوسيطة \textenglish{IrGenerator.cs} بالأنواع المؤكدة لإنتاج كود سليم ومحسن. \\
\end{tabularx}
\end{tcolorbox}

\begin{lstlisting}[language={[Sharp]C}, caption={مقتطف حقيقي من التحقق الدلالي وفحص مجالات الرؤية DoctorSemanticAnalyzer.cs}]
public override void VisitAssignment(AssignmentNode node)
{
    var symbol = _currentScope.Lookup(node.VariableName);
    if (symbol == null)
    {
        ReportError($"المعرف '{node.VariableName}' غير معرف في هذا النطاق.", node.Span);
        return;
    }
    if (symbol.IsReadOnly)
    {
        ReportError($"لا يمكن تعديل الثابت '{node.VariableName}' بعد تعيين قيمته الاولية.", node.Span);
        return;
    }

    LanguageType exprType = EvaluateType(node.ValueExpression);
    if (!IsTypeCompatible(symbol.Type, exprType))
    {
        ReportError($"تعارض في الانواع: لا يمكن اسناد قيمة من نوع '{exprType}' الى متغير من نوع '{symbol.Type}'.", node.Span);
    }
}
\end{lstlisting}

\newpage

% ------------------------------------------------------------------------------
\section{توليد الشفرة الوسيطة ثلاثية العناوين (IrGenerator.cs \& IrInstruction.cs)}

\begin{tcolorbox}[colback=white,colframe=vscodeBlue!80!black,arc=2mm,title=\textbf{بطاقة المواصفات البرمجية للشفرة الوسيطة (Three-Address Code Specification)}]
\begin{tabularx}{\linewidth}{p{4.2cm} X}
    \textbf{1. وظيفة الجزء من الكود} & تسطيح (Flattening) شجرة الإعراب الهرمية وتحويلها إلى متتالية خطية من أوامر العناوين الثلاثية (TAC)، وتوليد السجلات المؤقتة (temp\_0, temp\_1) وعناوين القفز الشرطي وغير الشرطي (Labels) لإدارة الحلقات والشروط. \\
    \midrule
    \textbf{2. المدخلات (Inputs)} & شجرة الإعراب المجردة \textenglish{ProgramNode} المؤكدة والمفحوصة دلالياً مع جدول الرموز. \\
    \midrule
    \textbf{3. المخرجات (Outputs)} & كائن \textenglish{IrProgram} يحتوي على قائمة الأوامر الوسيطة الخطية \textenglish{List\textless IrInstruction\textgreater}، خريطة حجز الذاكرة والمؤقتات \textenglish{StorageMap}، وجدول النصوص الثابتة \textenglish{StringLiterals}. \\
    \midrule
    \textbf{4. المتغيرات المستخدمة} & \textenglish{\_tempCounter} (عداد المؤقتات الخطية)، \textenglish{\_labelCounter} (عداد العلامات والعناوين)، \textenglish{\_instructions} (قائمة الأوامر الناتجة)، \textenglish{\_storageMap} (خريطة تسجيل أسماء المتغيرات وأنواعها). \\
    \midrule
    \textbf{5. هياكل البيانات المستخدمة} & سجل \textenglish{IrInstruction} (يحتوي على: Opcode, Left, Right, Result, TargetLabel, Span)، كلاس \textenglish{IrValue} (يمثل: Immediate, Variable, Temporary, StringLabel)، وتعداد \textenglish{IrOpcode} (أوامر الحساب، المقارنات، القفز، الطباعة، والقراءة). \\
    \midrule
    \textbf{6. الدوال أو الإجراءات المستخدمة} & \textenglish{Generate()} (بدء التوليد الوسيط)، \textenglish{GenerateExpression()} (تسطيح التعبيرات إلى مؤقتات)، \textenglish{NewTemporary()} (إنشاء سجل وسيط جديد)، \textenglish{NewLabel()} (إنشاء عنوان قفز فريد)، \textenglish{Emit()} (حقن تعليمة في البرنامج). \\
    \midrule
    \textbf{7. المكتبات المستخدمة} & \textenglish{System, System.Collections.Generic, BayanCompiler.Syntax, BayanCompiler.Intermediate}. \\
    \midrule
    \textbf{8. الملفات المرتبطة بالجزء} & \textenglish{IrProgram.cs, IrInstruction.cs, IrValue.cs, IrOpcode.cs, IrPrinter.cs}. \\
    \midrule
    \textbf{9. كيفية ارتباطه بباقي المشروع} & يمثل جسر العبور الفاصل بين الواجهة الأمامية للمترجم (Frontend) والواجهات الخلفية لتوليد الكود الهدف (Backend: MIPS, x86, CIL). \\
\end{tabularx}
\end{tcolorbox}

\begin{lstlisting}[language={[Sharp]C}, caption={مقتطف حقيقي من توليد اوامر الشفرة الوسيطة IrGenerator.cs}]
private IrValue GenerateBinaryExpression(BinaryExpressionNode binary)
{
    IrValue leftVal = GenerateExpression(binary.Left);
    IrValue rightVal = GenerateExpression(binary.Right);
    IrValue temp = NewTemporary(binary.InferredType);

    _program.Emit(new IrInstruction(
        IrOpcode.Binary,
        temp,
        leftVal,
        rightVal,
        binary.Operator,
        null,
        binary.Span
    ));
    return temp;
}
\end{lstlisting}

\newpage

% ------------------------------------------------------------------------------
\section{توليد لغة التجميع الهدفية لمعالجات x86 و MIPS (AssemblyGenerator.cs \& MipsCodeGenerator.cs)}

\begin{tcolorbox}[colback=white,colframe=vscodeBlue!80!black,arc=2mm,title=\textbf{بطاقة المواصفات البرمجية لتوليد كود التجميع (Target Assembly Generator Specification)}]
\begin{tabularx}{\linewidth}{p{4.2cm} X}
    \textbf{1. وظيفة الجزء من الكود} & تحويل تعليمات الشفرة الوسيطة (TAC) إلى شفرة تجميع هدفية حقيقية: إنتاج كود MIPS Assembly أصيل يوضح معمارية المعالجات (الأقسام \textenglish{.data} و \textenglish{.text}، والسجلات \textenglish{\$t0-\$t9}، وتعليمات الذاكرة \textenglish{lw} و \textenglish{sw} واستدعاءات النظام \textenglish{syscall})، وتوليد كود x86 32-bit. \\
    \midrule
    \textbf{2. المدخلات (Inputs)} & برنامج الشفرة الوسيطة ثلاثية العناوين \textenglish{IrProgram} ومسار مجلد الإخراج. \\
    \midrule
    \textbf{3. المخرجات (Outputs)} & ملف أسمبلي كامل \textenglish{output.asm} يحفظ في مجلد المخرجات ويعرض في التبويب المخصص بالمحرر. \\
    \midrule
    \textbf{4. المتغيرات المستخدمة} & \textenglish{builder} (لبناء النصوص المعمارية للأسمبلي)، \textenglish{fields} (جدول حقول الذاكرة والمواقع المحجوزة)، \textenglish{labels} (جدول علامات القفز)، \textenglish{allocatedRegisters} (خريطة تخصيص السجلات). \\
    \midrule
    \textbf{5. هياكل البيانات المستخدمة} & \textenglish{StringBuilder}، قواميس المعرفات \textenglish{Dictionary\textless string, string\textgreater}، مصفوفات حجز الذاكرة للبيانات الأولية والمركبة. \\
    \midrule
    \textbf{6. الدوال أو الإجراءات المستخدمة} & \textenglish{Generate()} (تنسيق توليد الملف)، \textenglish{EmitDataSection()} (بناء قسم حجز المتغيرات والسلاسل)، \textenglish{EmitTextSection()} (بناء قسم التعليمات التنفيذية)، \textenglish{TranslateInstruction()} (تحويل كل أمر TAC لأمر تجميع مقلوب). \\
    \midrule
    \textbf{7. المكتبات المستخدمة} & \textenglish{System, System.IO, System.Text, BayanCompiler.Intermediate}. \\
    \midrule
    \textbf{8. الملفات المرتبطة بالجزء} & \textenglish{MipsCodeGenerator.cs, AssemblyGenerator.cs, BackendTypeLayout.cs}. \\
    \midrule
    \textbf{9. كيفية ارتباطه بباقي المشروع} & يأخذ الشفرة الوسيطة من \textenglish{IrGenerator.cs}، وينتج الملف \textenglish{output.asm} الذي يمثل متطلب مرحلة الأسمبلي الرسمية في التقرير والمحرر. \\
\end{tabularx}
\end{tcolorbox}

\begin{lstlisting}[language={[Sharp]C}, caption={مقتطف حقيقي من توليد شفرة لغة التجميع MipsCodeGenerator.cs}]
private void EmitMipsInstruction(IrInstruction inst, StringBuilder sb)
{
    switch (inst.Opcode)
    {
        case IrOpcode.Assign:
            LoadValueToReg(inst.Left, "$t0", sb);
            sb.AppendLine($"    sw $t0, {GetMemoryLocation(inst.Result)}");
            break;
        case IrOpcode.Binary:
            LoadValueToReg(inst.Left, "$t0", sb);
            LoadValueToReg(inst.Right, "$t1", sb);
            string mipsOp = GetMipsOpcode(inst.Operator);
            sb.AppendLine($"    {mipsOp} $t2, $t0, $t1");
            sb.AppendLine($"    sw $t2, {GetMemoryLocation(inst.Result)}");
            break;
        case IrOpcode.PrintInt:
            LoadValueToReg(inst.Left, "$a0", sb);
            sb.AppendLine("    li $v0, 1");       // syscall 1: طباعة عدد صحيح
            sb.AppendLine("    syscall");
            break;
    }
}
\end{lstlisting}

\newpage

% ------------------------------------------------------------------------------
\section{توليد الملف التنفيذي الأصلي CIL و PE EXE (CilExecutableGenerator.cs \& ilasm.exe)}

\begin{tcolorbox}[colback=white,colframe=vscodeBlue!80!black,arc=2mm,title=\textbf{بطاقة المواصفات البرمجية لتوليد الملف التنفيذي (PE Executable Generator Specification)}]
\begin{tabularx}{\linewidth}{p{4.2cm} X}
    \textbf{1. وظيفة الجزء من الكود} & توليد كود لغة مايكروسوفت الوسيطة (CIL)، وضبط ترميز UTF-8 للمدخلات والمخرجات لحل مشكلة النصوص العربية، ثم استدعاء مجمع مايكروسوفت \textenglish{ilasm.exe} لتجميع الكود وإنتاج ملف تنفيذي أصلي \textenglish{output.exe} يعمل مباشرة في بيئة ويندوز دون محاكٍ. \\
    \midrule
    \textbf{2. المدخلات (Inputs)} & برنامج الشفرة الوسيطة \textenglish{IrProgram} ومسار مجلد الحفظ المختار. \\
    \midrule
    \textbf{3. المخرجات (Outputs)} & ملف الكود المصدري الوسيط \textenglish{output.il}، والملف التنفيذي النهائي المستقل \textenglish{output.exe} (Windows Portable Executable). \\
    \midrule
    \textbf{4. المتغيرات المستخدمة} & \textenglish{ilasmPath} (مسار أداة ilasm على نظام ويندوز)، \textenglish{ilText} (نص شفرة CIL المولدة)، \textenglish{locals} (خريطة مواقع المتغيرات المحلية في إطار المكدس)، \textenglish{startInfo} (معلمات إطلاق عملية المجمع). \\
    \midrule
    \textbf{5. هياكل البيانات المستخدمة} & \textenglish{ProcessStartInfo}، قواميس المتغيرات \textenglish{Dictionary\textless string, int\textgreater} لترقيم مواقع CIL Local Slots، وقواميس عنونة العلامات \textenglish{Dictionary\textless string, string\textgreater}. \\
    \midrule
    \textbf{6. الدوال أو الإجراءات المستخدمة} & \textenglish{Generate()} (تنسيق خطوات التوليد)، \textenglish{EmitMainHeader()} (بناء ترويسة دالة Main وضبط ترميز UTF8)، \textenglish{CompileIlToExe()} (إطلاق مجمع ilasm كعملية فرعية)، \textenglish{FindIlasmPath()} (البحث التلقائي عن مسار مجمع مايكروسوفت). \\
    \midrule
    \textbf{7. المكتبات المستخدمة} & \textenglish{System, System.Diagnostics, System.IO, System.Text}. \\
    \midrule
    \textbf{8. الملفات المرتبطة بالجزء} & \textenglish{CilExecutableGenerator.cs, AssemblyGenerator.cs, WindowsExecutableBuilder.cs}. \\
    \midrule
    \textbf{9. كيفية ارتباطه بباقي المشروع} & يمثل محطة الإنتاج الختامية للمترجم؛ يُنتج الملف التنفيذي الحقيقي الذي يقوم المحرر باستدعائه وتشغيله تفاعلياً داخل الطرفية المدمجة. \\
\end{tabularx}
\end{tcolorbox}

\begin{lstlisting}[language={[Sharp]C}, caption={مقتطف حقيقي من ضبط ترميز UTF-8 وتوليد الملف التنفيذي عبر ilasm}]
private void EmitMainHeader(StringBuilder builder, IrProgram program)
{
    builder.AppendLine("    .method public static void Main() cil managed");
    builder.AppendLine("    {");
    builder.AppendLine("        .entrypoint");
    builder.AppendLine("        .maxstack 32");
    builder.AppendLine();
    // فرض ترميز UTF-8 في دالة المدخل لضمان طباعة وقراءة النصوص العربية
    builder.AppendLine("        call class [mscorlib]System.Text.Encoding [mscorlib]System.Text.Encoding::get_UTF8()");
    builder.AppendLine("        call void [mscorlib]System.Console::set_OutputEncoding(class [mscorlib]System.Text.Encoding)");
    builder.AppendLine("        call class [mscorlib]System.Text.Encoding [mscorlib]System.Text.Encoding::get_UTF8()");
    builder.AppendLine("        call void [mscorlib]System.Console::set_InputEncoding(class [mscorlib]System.Text.Encoding)");
}
\end{lstlisting}

\newpage

% ------------------------------------------------------------------------------
\section{محرر لغة بيان وبيئة التطوير والطرفية التفاعلية (CompilerJourneyForm.cs \& BayanTerminalControl.cs)}

\begin{tcolorbox}[colback=white,colframe=vscodeBlue!80!black,arc=2mm,title=\textbf{بطاقة المواصفات البرمجية لمحرر اللغة والطرفية (Language Editor \& Terminal Specification)}]
\begin{tabularx}{\linewidth}{p{4.2cm} X}
    \textbf{1. وظيفة الجزء من الكود} & توفير بيئة تطوير برمجية متكاملة (IDE) مستقلة تحاكي واجهة Visual Studio Code، توفر التحرير متعدد التبويبات، تلوين الكلمات المفتاحية، شجرة العينات، فحص مخرجات مراحل المترجم التسع، استدعاء المترجم المستقل، وطرفية تفاعلية لتشغيل البرامج الحية مع استقبال المدخلات (Read). \\
    \midrule
    \textbf{2. المدخلات (Inputs)} & تفاعل المبرمج عبر لوحة المفاتيح والفأرة، ملفات الكود المصدري \textenglish{*.bayan}، ومدخلات الطرفية الحية أثناء التشغيل. \\
    \midrule
    \textbf{3. المخرجات (Outputs)} & شاشة تحرير عربية RTL متطورة، تبويبات عرض نتائج كل مرحلة (الرموز، الشجرة، جدول الرموز، TAC، الأسمبلي، الأخطاء)، والمخرجات الحية للبرامج المنفذة. \\
    \midrule
    \textbf{4. المتغيرات المستخدمة} & \textenglish{sourceEditor} (محرر الكود)، \textenglish{inspectorTabs} (فاحص المراحل التسع)، \textenglish{bottomTabs} (اللوحة السفلية)، \textenglish{terminalControl} (الطرفية التفاعلية)، \textenglish{compilerClient} (عميل الاتصال بالمترجم المستقل)، \textenglish{executableRuntimeClient} (مشغل الملف التنفيذي). \\
    \midrule
    \textbf{5. هياكل البيانات المستخدمة} & كلاس \textenglish{BayanDocument} (لحفظ بيانات الملف المفتوح وحالته)، مصفوفة الوثائق \textenglish{List\textless BayanDocument\textgreater}، حاويات التقسيم المرنة \textenglish{SplitContainer}، عناصر \textenglish{RichTextBox} للعرض النقي، وشجرة الاستكشاف \textenglish{TreeView}. \\
    \midrule
    \textbf{6. الدوال أو الإجراءات المستخدمة} & \textenglish{CompileSource()} (استدعاء عملية الترجمة عبر IPC)، \textenglish{RunExecutable()} (تشغيل الملف التنفيذي والربط بالطرفية)، \textenglish{BuildArabicVsCodeLayout()} (هندسة الواجهة والتقسيمات)، \textenglish{NavigateToStage()} (القفز الفوري للمرحلة المختارة)، \textenglish{ToggleTheme()} (التبديل بين الوضع الداكن والفاتح). \\
    \midrule
    \textbf{7. المكتبات المستخدمة} & \textenglish{System.Windows.Forms, System.Drawing, System.Diagnostics, BayanCompiler.Forms}. \\
    \midrule
    \textbf{8. الملفات المرتبطة بالجزء} & \textenglish{CompilerJourneyForm.cs, BayanTerminalControl.cs, BayanCodeEditor.cs, CliCompilationClient.cs, WindowsExecutableRuntimeClient.cs, IdeTheme.cs}. \\
    \midrule
    \textbf{9. كيفية ارتباطه بباقي المشروع} & يمثل الواجهة المرئية الجامعة للنظام بأكمله؛ يحقق معيار الفصل التام (Decoupling) باستدعاء المترجم المستقل كعملية خارجية واستقبال مخرجاته، ثم تشغيل الملف التنفيذي وإدارته داخل الطرفية. \\
\end{tabularx}
\end{tcolorbox}

\begin{lstlisting}[language={[Sharp]C}, caption={مقتطف حقيقي من استدعاء المترجم المستقل وعرض نتائج المراحل CompilerJourneyForm.cs}]
private async void CompileSource()
{
    SetBusy(true, "يجري تشغيل مترجم بيان المستقل والتحقق من القواعد...");
    try
    {
        // استدعاء المترجم المستقل عبر IPC
        CliCompilationRun run = await compilerClient.CompileAsync(sourceEditor.SourceCode);

        // توزيع مخرجات المراحل على تبويبات الفاحص
        txtTokens.Text = run.GetArtifact("tokens.txt");
        txtParseTree.Text = run.GetArtifact("parse-tree.txt");
        txtSymbolTable.Text = run.GetArtifact("symbol-table.txt");
        txtTac.Text = run.GetArtifact("three-address-code.txt");
        txtAssembly.Text = run.GetArtifact("output.asm");

        lastExecutablePath = run.Manifest?.WindowsExecutablePath ?? string.Empty;
        btnQuickRun.Enabled = run.Manifest?.Success == true && File.Exists(lastExecutablePath);
    }
    finally { SetBusy(false); }
}
\end{lstlisting}


\newpage

% ==============================================================================
% الصفحات الأخيرة: التوثيق البصري ونتائج تنفيذ المشروع
% ==============================================================================
\chapter{التوثيق البصري ونتائج الاختبار والتنفيذ}

يحتوي هذا الفصل على الصور التوثيقية الشاملة لمراحل تنفيذ المشروع، واجهات المحرر، ونتائج اختبار البرامج وفق متطلبات التقرير الأكاديمي.

\section{مخطط المعمارية ومسار الترجمة}
يوضح الشكل~(\ref{fig:arch}) المعمارية المنفصلة بين محرر الكود العربي ومحرك المترجم، وتدفق البيانات عبر المراحل التسع للترجمة.

\begin{figure}[H]
    \centering
    \includegraphics[width=0.95\textwidth]{images/arch_pipeline.png}
    \caption{مخطط معمارية النظام ومسار مراحل الترجمة التساعي للغة بيان}
    \label{fig:arch}
\end{figure}

\section{واجهة محرر بيان المستوحاة من VS Code}
يوضح الشكل~(\ref{fig:ide}) الواجهة الرئيسية للمحرر، وتظهر بالهوية البصرية العربية مع شريط النشاط اليميني، شجرة الأمثلة، محرر الكود، مفتش المراحل، والطرفية التفاعلية.

\begin{figure}[H]
    \centering
    \includegraphics[width=0.95\textwidth]{images/ide_overview.png}
    \caption{واجهة محرر لغة بيان العربية (Bayan Arabic VS Code Studio)}
    \label{fig:ide}
\end{figure}

\section{إدارة الملفات وقوائم الفتح والحفظ (File Management)}
يوضح الشكل~(\ref{fig:fileops}) قائمة «ملف» في المحرر، وتظهر خيارات فتح الملفات، الحفظ، الحفظ باسم، وإنشاء الملفات الجديدة مع دعم فتح ملفات متعددة في تبويبات مستقلة.

\begin{figure}[H]
    \centering
    \includegraphics[width=0.95\textwidth]{images/file_dialog.png}
    \caption{إدارة الملفات وقوائم الفتح والحفظ وتعدد التبويبات البرمجية}
    \label{fig:fileops}
\end{figure}

\newpage

\section{مخرجات التحليل المعجمي وجدول الرموز المعجمية}
يوضح الشكل~(\ref{fig:tokens}) جدول الرموز المعجمية الناتجة من تحليل كود البرنامج، مع تفصيل الفئة، القيمة المعجمية، ورقم السطر والعمود.

\begin{figure}[H]
    \centering
    \includegraphics[width=0.92\textwidth]{images/stage1_tokens.png}
    \caption{المرحلة الأولى: جدول تدفق الرموز المعجمية (Tokens Stream Table)}
    \label{fig:tokens}
\end{figure}

\section{شجرة الإعراب المجردة (AST)}
يوضح الشكل~(\ref{fig:ast}) الهيكل الهرمي لعقد شجرة الإعراب المجردة الناتجة من مطابقة قواعد $LL(1)$.

\begin{figure}[H]
    \centering
    \includegraphics[width=0.92\textwidth]{images/stage2_ast.png}
    \caption{المرحلتان الثانية والثالثة: شجرة الإعراب المجردة (Abstract Syntax Tree)}
    \label{fig:ast}
\end{figure}

\newpage

\section{جدول الرموز ومجالات الرؤية (Symbol Table)}
يوضح الشكل~(\ref{fig:symbols}) جدول الرموز وتفاصيل الأنواع ومجالات الرؤية (Scopes) وحالات التهيئة.

\begin{figure}[H]
    \centering
    \includegraphics[width=0.92\textwidth]{images/stage3_symbols.png}
    \caption{المرحلة الرابعة: تفاصيل جدول الرموز ومجالات الرؤية (Symbol Table Scopes)}
    \label{fig:symbols}
\end{figure}

\section{محرك التحليل الدلالي وفحص الأنواع}
يوضح الشكل~(\ref{fig:semantics}) عمليات التحقق الأربع التي يفرضها المحلل الدلالي للتأكد من سلامة البرنامج قبل التوليد.

\begin{figure}[H]
    \centering
    \includegraphics[width=0.92\textwidth]{images/stage4_semantics.png}
    \caption{المرحلة الرابعة: محرك التحليل الدلالي ومطابقة الأنواع ومجالات الرؤية}
    \label{fig:semantics}
\end{figure}

\newpage

\section{الشفرة الوسيطة ثلاثية العناوين (TAC)}
يوضح الشكل~(\ref{fig:tac}) التعليمات الخطية الوسيطة للشفرة ثلاثية العناوين وسجلاتها المؤقتة.

\begin{figure}[H]
    \centering
    \includegraphics[width=0.92\textwidth]{images/stage5_tac.png}
    \caption{المرحلة الخامسة: الشفرة الوسيطة ثلاثية العناوين (Three-Address Code)}
    \label{fig:tac}
\end{figure}

\section{توليد كود التجميع الهدفي لمعالجات x86}
يوضح الشكل~(\ref{fig:asm}) كود التجميع الهدفي المولد لمعالجات $x86$ متوافقاً مع مجمعات NASM/MASM.

\begin{figure}[H]
    \centering
    \includegraphics[width=0.92\textwidth]{images/stage6_assembly.png}
    \caption{المرحلة السادسة: كود التجميع الهدفي لمعالجات x86 32-bit Assembly}
    \label{fig:asm}
\end{figure}

\newpage

\section{توليد الملف التنفيذي الأصلي عبر ilasm}
يوضح الشكل~(\ref{fig:exe}) خط أنابيب توليد الملف التنفيذي الحقيقي $EXE$ وتأكيد تشغيله الذاتي.

\begin{figure}[H]
    \centering
    \includegraphics[width=0.92\textwidth]{images/stage7_executable.png}
    \caption{المرحلة السابعة: توليد الملف التنفيذي المستقل بنظام Windows PE عبر ilasm.exe}
    \label{fig:exe}
\end{figure}

\section{تشخيص الأخطاء والقفز لموضع الخطأ}
يوضح الشكل~(\ref{fig:errors}) قائمة الأخطاء النحوية والدلالية، وتوضيح ميزة القفز المباشر للسطر عند النقر المزدوج.

\begin{figure}[H]
    \centering
    \includegraphics[width=0.92\textwidth]{images/stage8_errors.png}
    \caption{المرحلة الثامنة: قائمة تشخيص الأخطاء وميزة القفز التلقائي للسطر والعمود}
    \label{fig:errors}
\end{figure}

\newpage

\section{تشغيل البرامج داخل الطرفية التفاعلية المدمجة}
يوضح الشكل~(\ref{fig:term}) مخرجات التنفيذ الحي لبرنامج مكتوب بلغة بيان داخل الطرفية المدمجة.

\begin{figure}[H]
    \centering
    \includegraphics[width=0.92\textwidth]{images/stage9_terminal_run.png}
    \caption{المرحلة التاسعة: تشغيل البرنامج الناتج وعرض مخرجاته في الطرفية التفاعلية}
    \label{fig:term}
\end{figure}

\newpage

% ==============================================================================
% مصفوفة أمثلة الاختبار الـ 18
% ==============================================================================
\chapter{مصفوفة برامج الاختبار والتحقق الرسمية (12 تطبيقاً شاملاً)}

وفقاً للمعيار الرابع من وثيقة متطلبات المشروع، تم تصميم 12 برنامج اختبار شاملاً ومتقدماً تغطي شتى قواعد وتراكيب لغة بيان المعتمدة من أستاذ المقرر:

\begin{table}[H]
    \centering
    \small
    \caption{جدول البرامج الاختبارية الشاملة لتغطية قواعد لغة بيان بالكامل}
    \begin{tabularx}{\textwidth}{l p{5.5cm} X c}
        \toprule
        \textbf{رقم} & \textbf{اسم ملف الاختبار} & \textbf{القاعدة اللغوية المغطاة والتطبيق} & \textbf{حالة التنفيذ} \\
        \midrule
        01 & \textenglish{01\_تعريف\_المتغيرات\_والثوابت\_والحساب.bayan} & المتغيرات، الثوابت، الحسابات الهندسية (المستطيل والدائرة)، باقي القسمة & ناجح (Exit 0) \\
        02 & \textenglish{02\_الجمل\_الشرطية\_وتقديرات\_الطلاب.bayan} & القرارات والشرطيات (إذا ... فإن ... وإلا)، المقارنات المنطقية والتقديرات & ناجح (Exit 0) \\
        03 & \textenglish{03\_حلقات\_التكرار\_والمجاميع\_والمضروب.bayan} & حلقات التكرار (طالما ... استمر، كرر ... إلى ... أضف، أعد ... حتى) & ناجح (Exit 0) \\
        04 & \textenglish{04\_الإجراءات\_وتمرير\_المعاملات.bayan} & الإجراءات الفرعية وتمرير المعاملات بالقيمة (ByVal) وبالمرجع (ByRef) & ناجح (Exit 0) \\
        05 & \textenglish{05\_القوائم\_والمصفوفات\_العددية.bayan} & القوائم (قائمة[...] من صحيح)، التعيين المفهرس، وحساب المجاميع & ناجح (Exit 0) \\
        06 & \textenglish{06\_السجلات\_وهياكل\_البيانات.bayan} & السجلات المركبة (سجل \{ ... \})، تمثيل النقاط، الوصول للحقول والرموز & ناجح (Exit 0) \\
        07 & \textenglish{07\_الحسابات\_العلمية\_والأعداد\_الحقيقية.bayan} & نوع حقيقي، حسابات الجاذبية والكتلة، مقارنات الأعداد العشرية & ناجح (Exit 0) \\
        08 & \textenglish{08\_الأس\_والقسمة\_الصحيحة\_المتقدمة.bayan} & العمليات الحسابية المتقدمة: معامل الأس (\^{})، معامل القسمة الصحيحة (\textbackslash) & ناجح (Exit 0) \\
        09 & \textenglish{09\_معالجة\_النصوص\_والمحارف.bayan} & السلاسل النصية (خيط رمزي)، المحارف (حرفي)، وطباعة الرسائل العربية & ناجح (Exit 0) \\
        10 & \textenglish{10\_تطبيق\_شامل\_متكامل\_لجميع\_القواعد.bayan} & تطبيق مالي متكامل يجمع: السجلات، القوائم، الإجراءات، الشروط، والتكرار & ناجح (Exit 0) \\
        11 & \textenglish{11\_اختبار\_الأخطاء\_النحوية\_المتعمدة.bayan} & التحقق من اكتشاف الأخطاء النحوية والقفز للسطر (Syntax Diagnostics) & كشف الخطأ بنجاح \\
        12 & \textenglish{12\_اختبار\_الأخطاء\_الدلالية\_المتعمدة.bayan} & التحقق من كشف تعارض الأنواع وحماية الثوابت (Semantic Diagnostics) & كشف الخطأ بنجاح \\
        \bottomrule
    \end{tabularx}
\end{table}

\chapter{الخاتمة والتوصيات}
تم بحمد الله وتوفيقه إنجاز مشروع مترجم ومحرر لغة **بيان** العربية محققاً 100\% من المتطلبات والشروط المحددة في وثيقة المشروع. أثبت النظام كفاءة عالية في:
\begin{enumerate}
    \item الفصل التام بين المترجم المستقل (CLI) وبيئة التطوير المرئية (VS Code IDE).
    \item توليد ملفات تنفيذية حقيقية ومستقلة بصيغة Windows PE عبر محول CIL ومجمّع ilasm.exe.
    \item توفير تجربة مستخدم عربية أصيلة متكاملة مع سهولة استكشاف المراحل ومعالجة الأخطاء.
\end{enumerate}

\end{document}
